% rosetta/preamble.tex — 数学↔物理 罗塞塔词典 导言（独立于主译本 preamble）
\documentclass[11pt,a4paper]{article}
\usepackage[UTF8, fontset=windows, heading=true]{ctex}
\usepackage{amsmath,amssymb,amsthm}
\usepackage{bm}
\usepackage{geometry}
\geometry{left=2.4cm,right=2.4cm,top=2.6cm,bottom=2.6cm}
\usepackage{xcolor}
\usepackage{graphicx}
\usepackage{pdflscape}
\usepackage{enumitem}
\usepackage{booktabs}
\usepackage{array}
\usepackage{longtable}
\usepackage{caption}
\usepackage{tcolorbox}
\tcbuselibrary{skins,breakable}
\usepackage[colorlinks=true,linkcolor=blue!45!black,urlcolor=blue!45!black]{hyperref}
\linespread{1.16}

% ---- 颜色 ----
\definecolor{mcol}{RGB}{28,78,140}    % 数学侧
\definecolor{pcol}{RGB}{150,74,10}    % 物理侧
\definecolor{gcol}{RGB}{105,105,105}

% ---- 侧标签 ----
\newcommand{\sidelabel}[2]{{\color{#1}\bfseries #2}\quad}
\newcommand{\mathdef}[1]{\par\noindent\sidelabel{mcol}{数学侧｜定义}#1\par}
\newcommand{\physdef}[1]{\par\noindent\sidelabel{pcol}{物理侧｜定义}#1\par}
\newcommand{\mathex}[1]{\par\noindent\sidelabel{mcol}{数学侧｜例子}#1\par}
\newcommand{\physex}[1]{\par\noindent\sidelabel{pcol}{物理侧｜例子}#1\par}
\newcommand{\bridge}[1]{\par\noindent\sidelabel{gcol}{对应}#1\par}
\newcommand{\intuit}[1]{\par\noindent\sidelabel{gcol}{直觉}#1\par}
\newcommand{\pitfall}[1]{\par\noindent\sidelabel{gcol}{易错}#1\par}

% ---- 词条环境：#1 中文名 #2 英文名 #3 ascii 标签 #4 原文锚点 ----
\newenvironment{word}[4]{%
  \subsection{#1（#2）}%
  \label{word:#3}%
  \noindent{\footnotesize\color{gcol}原文：#4}\par\smallskip
  \begingroup\setlength{\parskip}{3pt}%
}{%
  \endgroup\medskip\hrule height 0.2pt\medskip
}

% ---- 轨道标题 ----
\newcommand{\track}[2]{\section{#1}\noindent{\small\color{gcol}#2}\par\medskip}

% ---- 对照总表用的小字 ----
\newcommand{\ttfamilyx}[1]{{\ttfamily\small #1}}
